\section{Тестирование и бенчмарк}

Корректность программы проверялась на примере из условия и на автоматически сгенерированных тестах.
В бенчмарке результат каждого поиска сравнивается с позициями, найденными с помощью
\texttt{std::string::find}. Отдельно проверялись следующие ситуации:
\begin{itemize}
\item образец встречается один раз;
\item образец встречается несколько раз;
\item образец отсутствует в тексте;
\item большое число случайных образцов;
\item периодический текст с большим числом пересекающихся вхождений;
\item текст из одинаковых символов.
\end{itemize}

Для сравнения каждый образец также ищется прямым вызовом \texttt{std::string::find}.

\subsection{Образец встречается один раз}

В отдельном тесте использовалась случайная строка длиной $30\,000$ символов и небольшой набор
подстрок, каждая из которых встречается ровно один раз в тексте. Таким образом, суммарное число
совпадений равно нулю, а задача сводится к проверке отрицательных и единичных случаев поиска.

\begin{center}
\small
\begin{tabular}{|l|r|r|r|r|r|}
\hline
\rowcolor{lightgray}
Тест & Длина текста & Образцов & Вхождений & СД search, \(\mu s\) & std::find, \(\mu s\) \\
\hline
\texttt{unique\_match} & 30\,000 & 10 & 0 & 3 & 115 \\
\hline
\end{tabular}
\end{center}

На данном тесте суффиксное дерево демонстрирует существенно более высокую производительность
по сравнению с последовательным поиском, поскольку после построения структуры каждый запрос
обрабатывается без повторного сканирования всего текста. В то же время \texttt{std::string::find}
выполняет линейный проход по строке для каждого образца, что приводит к заметному росту времени
при увеличении числа запросов.

\subsection{Случайные тексты}

В первой серии текст и часть образцов генерировались случайно. Половина образцов бралась как реальные
подстроки текста, остальные образцы генерировались независимо. В последнем столбце указано отношение
\texttt{std::find / SuffixTree search}; значение больше 1 означает, что поиск по построенному дереву
быстрее прямого повторного поиска по строке.

\begin{center}
\small
\begin{tabular}{|l|r|r|r|r|r|}
\hline
\rowcolor{lightgray}
Тест & Длина текста & Образцов & Build, \(\mu s\) & СД search, \(\mu s\) & std::find / СД \\
\hline
\texttt{random\_1000} & 1\,000 & 100 & 300 & 26 & \(5.308\times\) \\
\hline
\texttt{random\_10000} & 10\,000 & 500 & 3564 & 127 & \(38.709\times\) \\
\hline
\texttt{random\_30000} & 30\,000 & 1000 & 9641 & 258 & \(53.453\times\) \\
\hline
\end{tabular}
\end{center}

На случайных данных время построения дерева растёт вместе с длиной текста, а время обработки множества
образцов после построения остаётся существенно меньше времени построения. Суффиксное дерево показывает
существенное ускорение относительно повторного \texttt{std::string::find}, поскольку каждый запрос
обрабатывается за время, не зависящее от длины текста после построения структуры.

\subsection{Много вхождений}

Во второй серии проверялись тексты с большим числом ответов. В таких тестах существенная часть времени
тратится не на прохождение по образцу, а на обход поддерева, сбор всех листьев и сортировку позиций.

\begin{center}
\small
\begin{tabular}{|l|r|r|r|r|r|}
\hline
\rowcolor{lightgray}
Тест & Длина текста & Образцов & Вхождений & СД search, \(\mu s\) & std::find, \(\mu s\) \\
\hline
\texttt{periodic\_30000} & 30\,000 & 500 & 1\,114\,016 & 90476 & 32128 \\
\hline
\texttt{same\_20000} & 20\,000 & 20 & 399\,810 & 15162 & 4489 \\
\hline
\end{tabular}
\end{center}

На периодическом тексте прямой поиск оказался быстрее: отношение
\texttt{std::find / SuffixTree search} составляет \(0.355\times\), то есть
\texttt{std::string::find} работает примерно в \(2.8\times\) быстрее.

На строке из одинаковых символов стандартный поиск также быстрее:
отношение равно \(0.296\times\), что соответствует преимуществу
\texttt{std::string::find} примерно в \(3.4\times\). Это объясняется тем, что
в сценариях с огромным числом совпадений стоимость обхода дерева и обработки
результатов становится доминирующей по сравнению с линейным сканированием.
\pagebreak